\documentclass[11pt]{article}

\usepackage[margin=1in]{geometry}
\usepackage[T1]{fontenc}
\usepackage[utf8]{inputenc}
\usepackage{lmodern}
\usepackage{amsmath,amssymb}
\usepackage{booktabs}
\usepackage{array}
\usepackage{graphicx}
\usepackage{subcaption}
\usepackage{placeins}
\usepackage{hyperref}

\hypersetup{colorlinks=true,linkcolor=blue,urlcolor=blue,citecolor=blue}

\newcommand{\FigDir}{../../figures}
\newcommand{\Aeff}{\mathcal A_{\mathrm{eff}}}
\newcommand{\Liq}{\mathcal L}
\newcommand{\Fresp}{\mathcal F}
\newcommand{\sigloc}{\sigma_{\mathrm{loc}}}

\title{Supplementary Materials:\\Reaction-Boundary Volatility Surface Simulations}
\author{}
\date{}

\begin{document}
\maketitle

\section{Short Background}

This supplementary note describes a simulation study of volatility surfaces
generated from the reaction-boundary closure. The purpose is not calibration.
Instead, the figures show how structural order-book channels can produce
stock-price and maturity dependence in a projected local-volatility surface.
The simulations are intended as a visual bridge between the theoretical
closure and the kinds of surfaces one would later compare with calibrated
local-volatility data.

\section{Reaction-Boundary Volatility Closure}

The operational variance closure used in the simulation is
\begin{equation}
 a_u^{(\Delta)}(S,u)
 \simeq
 \left[
 \frac{\Aeff(S,u)}
      {\Liq(S,u)^2}
 \right]
 \Delta^{-\gamma(S,u)}
 \Fresp_{\gamma(S,u)}\!\left(\nu(S,u)\Delta\right).
 \label{eq:operational-closure}
\end{equation}
Here \(S\) denotes the stock-price coordinate, \(u\) is operational time,
\(\Aeff\) is an effective signed-forcing scale, \(\Liq\) is the local
liquidity slope, \(\Delta\) is the operational coarse-graining scale,
\(\gamma\) is the long-memory exponent, and
\(z=\nu(S,u)\Delta\) is the dimensionless resilience-response argument.

Under a deterministic activity clock, the projected local variance is written
as
\begin{equation}
 \sigloc^2(S,t)
 =
 \alpha(t)\,
 \Xi(S,U(t);\Delta),
 \label{eq:projected-local-variance}
\end{equation}
where \(\alpha(t)\) maps operational-time variance into calendar time. In the
simulations, \(\Xi\) is evaluated using the right-hand side of
Eq.~\eqref{eq:operational-closure}. The plotted quantity is projected local
volatility,
\begin{equation}
 \sigloc(m,T)=\sqrt{\sigloc^2(m,T)},
 \label{eq:plotted-volatility}
\end{equation}
on a grid in log-moneyness \(m=\ln(S/S_0)\) and maturity \(T\).

\section{Response-Regime Approximations}

Both asymptotic response regimes are simulated. The low-\(z\) branch uses
\begin{equation}
 \Fresp_\gamma(z)
 \sim
 \frac{\tan(\pi\gamma/2)}{2(1-\gamma)},
 \qquad z\ll1,
 \label{eq:low-z-response}
\end{equation}
which removes the leading-order dependence on \(z\) when \(\gamma\) is held
fixed. The high-\(z\) branch uses
\begin{equation}
 \Fresp_\gamma(z)
 \sim
 K_\gamma^{(0)}z^{\gamma-1},
 \qquad z\gg1,
 \label{eq:high-z-response}
\end{equation}
where
\begin{equation}
 K_\gamma^{(0)}
 =
 \frac{
 \Gamma(1-\gamma)\sin(\pi\gamma/2)
 \Gamma(\gamma/2)\Gamma((1-\gamma)/2)
 }{2\pi^{3/2}}.
 \label{eq:k-gamma}
\end{equation}
Because \(0<\gamma<1\), the exponent \(\gamma-1\) is negative in the
high-\(z\) branch. Smaller \(z\), corresponding here to weaker or slower
resilience at fixed \(\Delta\), therefore raises the response multiplier in
that branch.

\begin{table}[ht]
\centering
\small
\caption{Response-function regimes used in the simulations.}
\label{tab:response-regimes}
\begin{tabular}{p{0.19\linewidth}p{0.23\linewidth}p{0.37\linewidth}p{0.11\linewidth}}
\toprule
Regime & Condition & Approximation & Depends on \(z\) \\
\midrule
Low \(z\) & \(z=\nu\Delta\ll1\) &
\(\tan(\pi\gamma/2)/(2(1-\gamma))\) & No \\
High \(z\) & \(z=\nu\Delta\gg1\) &
\(K_\gamma^{(0)}z^{\gamma-1}\) & Yes \\
\bottomrule
\end{tabular}
\end{table}

\section{Structural Simulation Profiles}

The stock-price coordinate is represented by log-moneyness,
\begin{equation}
 m=\ln(S/S_0).
\end{equation}
Stock dependence can enter the closure through \(\Liq(S)\), \(\Aeff(S)\),
\(\nu(S)\), \(\gamma(S)\), and possibly the activity clock. The first
simulation holds \(\gamma\) constant and uses a simple maturity-dependent
activity profile,
\begin{equation}
 \alpha(T)
 =
 \alpha_0\left(1+a_{\mathrm{short}}\exp(-T/\tau_\alpha)\right).
 \label{eq:activity-profile}
\end{equation}
This keeps the source of \(m\)-dependence visible in each profile.

\begin{table}[ht]
\centering
\small
\caption{Structural profiles used to generate stock dependence in the simulated surfaces.}
\label{tab:profile-definitions}
\begin{tabular}{p{0.22\linewidth}p{0.26\linewidth}p{0.22\linewidth}p{0.20\linewidth}}
\toprule
Profile & Varies with \(m\) & Held fixed & Purpose \\
\midrule
Liquidity-only & \(\Liq(m)\) &
\(\Aeff,z,\gamma,\alpha\) &
Isolates the inverse-square liquidity channel. \\
Resilience-response-only & \(z(m)=\nu(m)\Delta\) &
\(\Aeff,\Liq,\gamma,\alpha\) &
Isolates the response channel. \\
Combined structural & \(\Liq(m),z(m),\Aeff(m)\) &
\(\gamma\), simple \(\alpha(T)\) &
Combines stylised liquidity, resilience, and stress channels. \\
\bottomrule
\end{tabular}
\end{table}

The common functional forms are
\begin{align}
 \Liq(m)
 &=
 L_0\max\!\left(L_{\mathrm{floor}},
 1+L_{\mathrm{skew}}m+L_{\mathrm{smile}}m^2\right),
 \label{eq:liquidity-profile}\\
 z(m)
 &=
 z_{\mathrm{regime}}\exp(\nu_{\mathrm{skew}}m),
 \label{eq:z-profile}\\
 \Aeff(m)
 &=
 A_{\mathrm{eff},0}\exp(-A_{\mathrm{stress}}m).
 \label{eq:forcing-profile}
\end{align}
Positive \(L_{\mathrm{skew}}\) lowers liquidity for negative \(m\), while
positive \(L_{\mathrm{smile}}\) changes liquidity away from the money.
Positive \(\nu_{\mathrm{skew}}\) makes \(z(m)\) smaller on the downside.
In the combined profile, \(A_{\mathrm{stress}}\) raises the effective forcing
scale for negative \(m\). These choices are phenomenological simulation
controls, not fitted estimates.

\section{Numerical Implementation}

The surfaces are evaluated on a rectangular grid in \((m,T)\). For each grid
point, the selected structural profile determines \(\Liq(m)\), \(z(m)\), and
\(\Aeff(m)\); the selected response branch determines \(\Fresp_\gamma(z)\);
and Eqs.~\eqref{eq:operational-closure}--\eqref{eq:plotted-volatility}
determine the plotted volatility. The figures were generated with Python,
NumPy, and Matplotlib. Surface figures use boxed three-dimensional axes, and
the paired contour figures show black contour lines of the same surface.

\section{Relation To Calibrated Local Volatility}

The plotted object is best interpreted as a projected structural volatility
surface. It is local-volatility-like because it assigns a volatility level to
a state coordinate and maturity, but it is not a calibrated Dupire local
volatility surface. A calibrated local-volatility surface is usually inferred
from option prices under a risk-neutral pricing measure and must be consistent
with the absence of static arbitrage in the associated option surface. The
present construction instead starts from order-book structural quantities and
projects them into a surface morphology.

This distinction matters for interpretation. The simulation is designed to
ask whether plausible order-book channels can generate volatility-surface
shapes that resemble known market phenomenology, especially equity-like
downside skew. It does not claim that the resulting surface is directly
tradeable, arbitrage-free, or equivalent to a market-calibrated local
volatility input.

For comparison with option-calibrated surfaces, the state coordinate \(m\)
should be read as a log-moneyness coordinate. In a later empirical comparison,
one would normally align this with log-strike or log-forward moneyness, for
example \(m=\ln(K/S_0)\) or \(m=\ln(K/F_0(T))\), depending on the calibration
convention. The current notation \(m=\ln(S/S_0)\) keeps the connection to the
structural stock-price coordinate visible, while the plotted surface should
be compared only qualitatively with calibrated local volatility.

Within that qualitative interpretation, the combined structural high-\(z\)
case is the most market-like of the first simulations. It combines thinner
downside liquidity, weaker downside resilience through \(z(m)=\nu(m)\Delta\),
and stronger downside forcing. The low-\(z\) resilience-response-only case is
mainly diagnostic: at fixed \(\gamma\), the leading low-\(z\) approximation
does not depend on \(z\), so varying \(z(m)\) alone cannot generate leading
stock dependence in that branch.

\section{Parameter Choices}

The baseline values are dimensionless simulation controls. They set a
reasonable scale for visual comparison across profiles and regimes, but they
should not be read as empirical estimates.

\begin{table}[ht]
\centering
\small
\caption{Dimensionless baseline values and simulation ranges.}
\label{tab:parameter-ranges}
\begin{tabular}{p{0.18\linewidth}p{0.13\linewidth}p{0.13\linewidth}p{0.13\linewidth}p{0.31\linewidth}}
\toprule
Parameter & Baseline & Minimum & Maximum & Role \\
\midrule
\(\gamma\) & 0.55 & 0.30 & 0.80 & Long-memory exponent; avoids endpoints 0 and 1. \\
\(\Delta\) & 1.0 & 0.5 & 2.0 & Operational coarse-graining scale. \\
\(A_{\mathrm{eff},0}\) & 1.0 & 0.5 & 2.0 & Effective signed-forcing intensity scale. \\
\(L_0\) & 1.0 & 0.75 & 1.5 & Baseline liquidity slope. \\
\(z_{\mathrm{low}}\) & 0.05 & 0.01 & 0.10 & Low-\(z\) response-regime marker. \\
\(z_{\mathrm{high},0}\) & 10.0 & 5.0 & 25.0 & High-\(z\) response-regime marker. \\
\(\alpha_0\) & 1.0 & 0.5 & 2.0 & Baseline deterministic activity. \\
\(a_{\mathrm{short}}\) & 0.25 & 0.0 & 0.75 & Short-maturity activity lift. \\
\(\tau_\alpha\) & 0.50 & 0.25 & 1.0 & Activity decay scale in years. \\
\(L_{\mathrm{skew}}\) & 0.80 & 0.0 & 1.5 & Lower downside liquidity when positive. \\
\(L_{\mathrm{smile}}\) & 1.50 & 0.0 & 3.0 & Liquidity variation away from the center. \\
\(\nu_{\mathrm{skew}}\) & 1.00 & 0.0 & 2.0 & Weaker downside resilience when positive. \\
\(A_{\mathrm{stress}}\) & 0.75 & 0.0 & 1.5 & Stronger downside forcing in the combined profile. \\
\(L_{\mathrm{floor}}\) & 0.20 & 0.10 & 0.50 & Positive floor for the liquidity profile. \\
\bottomrule
\end{tabular}
\end{table}

\FloatBarrier

\section{Simulated Surfaces}

Each figure pair has the same interpretation. The left panel is the boxed
surface view of projected local volatility \(\sigloc(m,T)\). The right panel
is the black contour-line projection of the same surface. The horizontal
state coordinate is \(m=\ln(S/S_0)\), the maturity coordinate is \(T\) in
years, and the vertical surface coordinate is projected local volatility in
dimensionless simulation units.

\subsection{Liquidity-Only Profile}

\begin{figure}[ht]
 \centering
 \begin{subfigure}{0.48\linewidth}
  \centering
  \includegraphics[width=\linewidth]{\FigDir/mpl-surface3d-liquidity-only-low-z-v3.pdf}
  \caption{Boxed surface view.}
 \end{subfigure}
 \hfill
 \begin{subfigure}{0.48\linewidth}
  \centering
  \includegraphics[width=\linewidth]{\FigDir/mpl-contour-lines-liquidity-only-low-z-v3.pdf}
  \caption{Black contour-line projection.}
 \end{subfigure}
 \caption{Liquidity-only simulation in the low-\(z\) response regime. The only structural source of stock dependence is the liquidity profile \(\Liq(m)\), which enters the closure through the inverse-square factor \(1/\Liq(m)^2\). Since the low-\(z\) response approximation is independent of \(z\) at fixed \(\gamma\), the response channel does not contribute leading-order stock dependence in this panel.}
 \label{fig:liquidity-low-z}
\end{figure}

\begin{figure}[ht]
 \centering
 \begin{subfigure}{0.48\linewidth}
  \centering
  \includegraphics[width=\linewidth]{\FigDir/mpl-surface3d-liquidity-only-high-z-v3.pdf}
  \caption{Boxed surface view.}
 \end{subfigure}
 \hfill
 \begin{subfigure}{0.48\linewidth}
  \centering
  \includegraphics[width=\linewidth]{\FigDir/mpl-contour-lines-liquidity-only-high-z-v3.pdf}
  \caption{Black contour-line projection.}
 \end{subfigure}
 \caption{Liquidity-only simulation in the high-\(z\) response regime. The liquidity profile remains the only source of stock dependence because \(z\) is held constant in this profile. The comparison with Fig.~\ref{fig:liquidity-low-z} isolates the change in response-regime scale while preserving the same liquidity-driven surface shape.}
 \label{fig:liquidity-high-z}
\end{figure}

\FloatBarrier

\subsection{Resilience-Response-Only Profile}

\begin{figure}[ht]
 \centering
 \begin{subfigure}{0.48\linewidth}
  \centering
  \includegraphics[width=\linewidth]{\FigDir/mpl-surface3d-resilience-response-only-low-z-v3.pdf}
  \caption{Boxed surface view.}
 \end{subfigure}
 \hfill
 \begin{subfigure}{0.48\linewidth}
  \centering
  \includegraphics[width=\linewidth]{\FigDir/mpl-contour-lines-resilience-response-only-low-z-v3.pdf}
  \caption{Black contour-line projection.}
 \end{subfigure}
 \caption{Resilience-response-only simulation in the low-\(z\) response regime. The profile varies \(z(m)=\nu(m)\Delta\), but the leading low-\(z\) response approximation in Eq.~\eqref{eq:low-z-response} does not depend on \(z\) at fixed \(\gamma\). The figure is therefore a diagnostic check: it shows the maturity activity effect without an additional leading-order response-driven stock dependence.}
 \label{fig:resilience-low-z}
\end{figure}

\begin{figure}[ht]
 \centering
 \begin{subfigure}{0.48\linewidth}
  \centering
  \includegraphics[width=\linewidth]{\FigDir/mpl-surface3d-resilience-response-only-high-z-v3.pdf}
  \caption{Boxed surface view.}
 \end{subfigure}
 \hfill
 \begin{subfigure}{0.48\linewidth}
  \centering
  \includegraphics[width=\linewidth]{\FigDir/mpl-contour-lines-resilience-response-only-high-z-v3.pdf}
  \caption{Black contour-line projection.}
 \end{subfigure}
 \caption{Resilience-response-only simulation in the high-\(z\) response regime. Here \(z(m)\) enters directly through \(K_\gamma^{(0)}z^{\gamma-1}\). With \(\gamma<1\), lower downside \(z(m)\) increases the response multiplier, producing stock dependence even though liquidity and forcing are held fixed.}
 \label{fig:resilience-high-z}
\end{figure}

\FloatBarrier

\subsection{Combined Structural Profile}

\begin{figure}[ht]
 \centering
 \begin{subfigure}{0.48\linewidth}
  \centering
  \includegraphics[width=\linewidth]{\FigDir/mpl-surface3d-combined-structural-low-z-v3.pdf}
  \caption{Boxed surface view.}
 \end{subfigure}
 \hfill
 \begin{subfigure}{0.48\linewidth}
  \centering
  \includegraphics[width=\linewidth]{\FigDir/mpl-contour-lines-combined-structural-low-z-v3.pdf}
  \caption{Black contour-line projection.}
 \end{subfigure}
 \caption{Combined structural simulation in the low-\(z\) response regime. The surface combines liquidity variation and downside forcing while the leading response approximation remains independent of \(z\) at fixed \(\gamma\). The resulting shape should be interpreted as a stylised liquidity-and-forcing effect rather than as a calibrated volatility surface.}
 \label{fig:combined-low-z}
\end{figure}

\begin{figure}[ht]
 \centering
 \begin{subfigure}{0.48\linewidth}
  \centering
  \includegraphics[width=\linewidth]{\FigDir/mpl-surface3d-combined-structural-high-z-v3.pdf}
  \caption{Boxed surface view.}
 \end{subfigure}
 \hfill
 \begin{subfigure}{0.48\linewidth}
  \centering
  \includegraphics[width=\linewidth]{\FigDir/mpl-contour-lines-combined-structural-high-z-v3.pdf}
  \caption{Black contour-line projection.}
 \end{subfigure}
 \caption{Combined structural simulation in the high-\(z\) response regime. This profile combines the inverse-square liquidity channel, the response multiplier \(K_\gamma^{(0)}z^{\gamma-1}\), and a stylised downside forcing profile. The figure is the closest of the six simulations to the intended phenomenological comparison with a calibrated volatility surface, while remaining a mechanism simulation rather than an empirical fit.}
 \label{fig:combined-high-z}
\end{figure}

\FloatBarrier

\section{Robustness Checks}

The numerical checks verify that the simulated surfaces are finite and
positive on the grid, that the liquidity floor prevents invalid liquidity
values, and that the response-regime cases behave consistently with their
asymptotic formulas. The checks also verify the qualitative distinction
between the low-\(z\) and high-\(z\) response channels: when only \(z(m)\)
varies, the low-\(z\) branch has no leading-order \(m\)-dependence at fixed
\(\gamma\), while the high-\(z\) branch does.

\section{Limitations And Interpretation}

The simulations are structural illustrations. They do not estimate
\(\Liq(S)\), \(\nu(S)\), \(\Aeff(S)\), \(\gamma(S)\), or \(\alpha(T)\) from
market data. The stock dependence of the plotted surfaces is selected through
declared phenomenological profiles: inverse-square liquidity effects,
response dependence through \(z(S)=\nu(S)\Delta\), and a stylised downside
forcing scale. The role of the figures is therefore to make the theory
visually testable and to clarify which structural channels can plausibly
produce surface morphology before any empirical calibration is attempted.

\section*{Use of Generative AI}

Generative AI tools were used as computational and editorial assistance in
preparing this supplementary note, including drafting explanatory text,
organizing simulation code, and checking figure and table presentation. The
theoretical interpretation, modelling assumptions, parameter choices, code
outputs, and final scientific claims were reviewed by the author(s), who
remain responsible for the content. No empirical market data are used in
these simulations; the figures are generated from the stated equations and
declared parameter profiles.

\end{document}
